\phantomsection
\addcontentsline{toc}{section}{Décharge}
\thispagestyle{plain}

\vspace*{3cm}

\begin{center}
\begin{tikzpicture}
  % Carte principale : coins arrondis, dégradé doux, ombre portée
  \node[
    draw=ensaeblue,
    line width=1pt,
    rounded corners=10pt,
    top color=white,
    bottom color=ensaeblue!6,
    blur shadow={shadow blur steps=6, shadow xshift=1.5pt,
                 shadow yshift=-2.5pt, shadow opacity=45},
    inner xsep=32pt,
    inner ysep=36pt,
    text width=0.72\linewidth,
    align=justify,
    font=\itshape
  ] (box) {
    \vspace{6pt}%
    L'École nationale de la statistique et de l'analyse économique Pierre
    Ndiaye (ENSAE) et l'Agence nationale de la statistique et de la
    démographie (ANSD) n'entendent donner aucune approbation ni improbation
    aux idées et opinions émises dans ce mémoire. Ces idées et opinions
    doivent être considérées comme propres à leur auteur.
  };

  % Filet fin intérieur (effet double encadrement)
  \draw[ensaecyan, line width=0.4pt, rounded corners=7pt]
    ([xshift=6pt,yshift=-6pt]box.north west)
    rectangle
    ([xshift=-6pt,yshift=6pt]box.south east);

  % Filet d'accent bleu sur le bord gauche
  \draw[ensaeblue, line width=3pt, rounded corners=10pt]
    ([xshift=1pt,yshift=-14pt]box.north west) --
    ([xshift=1pt,yshift=14pt]box.south west);

  % Pastille de titre « Décharge » posée sur le bord supérieur
  \node[
    fill=ensaeblue,
    text=white,
    rounded corners=4pt,
    inner xsep=16pt,
    inner ysep=6pt,
    font=\normalfont\scshape\large,
    anchor=center
  ] at (box.north) {Décharge};

  % Guillemets en filigrane
  \node[anchor=north west, text=ensaeblue, opacity=0.22,
        font=\fontsize{50}{50}\selectfont\bfseries]
    at ([xshift=10pt,yshift=6pt]box.north west) {\guillemotleft};
  \node[anchor=south east, text=ensaeblue, opacity=0.22,
        font=\fontsize{50}{50}\selectfont\bfseries]
    at ([xshift=-10pt,yshift=-8pt]box.south east) {\guillemotright};
\end{tikzpicture}
\end{center}
